\subsection{Dolma、DCLM、Nemotron-CC：从多样性到可比较过滤}

Dolma 进一步把 Reddit（来自 Pushshift 项目，2005--2023 年的帖子和评论）、学术论文（来自 Semantic Scholar 的 4000 万篇论文，项目名 PeS2o）、Wikipedia、C4、Project Gutenberg 等来源组合起来，强调多样性。其 Common Crawl 处理使用 fastText 做语言识别、Gopher + C4 规则过滤（有意避免模型过滤）、Jigsaw 分类器做毒性过滤、Bloom filter 做去重。总量 3T tokens。

DataComp-LM (DCLM) 则试图把数据处理算法变成标准化基准。它先把 Common Crawl 处理成 DCLM-pool（240T tokens），然后用质量分类器过滤。分类器的训练非常值得注意：正例（200K 条）来自 OpenHermes-2.5（GPT-4 生成的指令数据）和 ELI5（Reddit 的 Explain Like I'm 5 子版块），负例（200K 条）来自 RefinedWeb 的随机样本。用 fastText 训练后，在整个 DCLM-pool 上运行，过滤后得到 3.8T tokens。实验表明，这种分类器过滤显著优于其他过滤方法。

Nemotron-CC 走得更激进：针对 FineWebEdu 和 DCLM 过滤太激进（删除 90\% 以上数据）的问题，它用分类器集成（Nemotron-340B-instruct 打分蒸馏出的快速模型 + DCLM 分类器）做质量评估，同时对低质量数据用语言模型做重写（rephrase），对高质量数据用语言模型生成任务（QA 对、关键信息提取等）。最终产出 6.3T tokens（其中高质量子集 1.1T）。

\begin{table}[H]
\centering
\begin{tabular}{p{0.16\textwidth}p{0.28\textwidth}p{0.27\textwidth}p{0.23\textwidth}}
\toprule
\textbf{项目} & \textbf{策略} & \textbf{优点} & \textbf{潜在代价} \\
\midrule
Dolma & 多源混合 & 覆盖面广，便于做系统比较 & 来源繁杂，清洗标准必须很清楚 \\
DCLM & 标准化数据基准 & 让“数据处理”变成可比较问题 & 容易把基准本身当成最终目标 \\
Nemotron-CC & 过滤 + 重写 + 加权 & 能把高质量样本放大 & 生成式重写会引入模型偏差 \\
\bottomrule
\end{tabular}
\caption{从多源混合到可比较过滤，再到生成式增强，每一步都在牺牲不同东西}
\end{table}

讲者对这类工作的关注点并不只是“谁更大”，而是“谁更能说明数据方法本身”。一旦你开始把 filtering、ranking、rewriting、mixing 分成可研究的子问题，数据处理就不再只是前处理脚本，而是模型能力的重要组成部分。

\begin{table}[H]
\centering
\begin{tabular}{p{0.19\textwidth}p{0.29\textwidth}p{0.39\textwidth}}
\toprule
\textbf{数据集} & \textbf{过滤方式} & \textbf{讲者强调的取向} \\
\midrule
CCNet & Wikipedia 风格的 n-gram 打分 & 质量优先，偏保守 \\
C4 & 大量规则 + 英文过滤 & 简单透明，覆盖较广 \\
The Pile & 人工策划的多域混合 & 多样性优先 \\
LLaMA / RedPajama & 多源混合 + 清晰整理 & 开源可复现 \\
FineWeb & 更细致的去重、语言识别、PII 清理 & 现代网页数据工程 \\
DCLM & 数据竞赛基准 + 模型式过滤 & 方法可比较 \\
Nemotron-CC & 过滤 + 生成式重写 + 质量加权 & 在高质量与规模之间折中 \\
\bottomrule
\end{tabular}
\caption{不同代际的数据集各自偏好的质量与规模平衡点}
\end{table}

